% Generated draft: requires mathematical and source review.
\documentclass[11pt]{article}
\usepackage[margin=25mm]{geometry}
\usepackage{fontspec}
\setmainfont{Latin Modern Roman}
\usepackage{amsmath,amssymb,mathtools,unicode-math}
\setmathfont{Latin Modern Math}
\usepackage{enumitem,needspace,xurl,hyperref}
\setlength{\parindent}{0pt}
\setlength{\parskip}{5pt}
\setlength{\emergencystretch}{3em}
\newcommand{\N}{\mathbb{N}}
\newcommand{\Z}{\mathbb{Z}}
\newcommand{\Q}{\mathbb{Q}}
\newcommand{\R}{\mathbb{R}}
\newcommand{\C}{\mathbb{C}}
\newcommand{\F}{\mathbb{F}}
\newcommand{\A}{\mathbb{A}}
\newcommand{\PP}{\mathbb{P}}
\newcommand{\E}{\mathbb{E}}
\newcommand{\eps}{\varepsilon}
\newcommand{\dd}{\,\mathrm{d}}
\begin{document}
\section*{UnsolvedMath ID: 801098}
\textbf{Source problem code:} catalogue entry
\subsection{Definitions and mathematical statement}
Let \(K\) be a number field, fix an algebraic closure \(\overline K\), let \(B\geq 1\) be a degree bound, and let \(d\geq 2\) be an integer. For \(P=(x,y)\in(\overline K)^2\), set \(K(P)=K(x,y)\). Write \(\mathcal H_d(K)\) for the family referred to in the supplied record as generalized H\'{e}non maps of degree \(d\) over \(K\). For \(f\in\mathcal H_d(K)\), let \(h(f)\) denote the height of its tuple of coefficients, and retain the notation \(\hat h_f(P)\) used in the source.

The supplied excerpt does not define the precise family \(\mathcal H_d(K)\), the function \(\hat h_f\), or the coefficient-height conventions. Consequently, the following statement retains these source-dependent objects and requires their definitions to be checked against the cited source.

The conjecture asserts that there exist constants \(\epsilon>0\), \(C\in\R\), and \(N\in\Z_{\geq 0}\), depending only on \(K,B,d\), such that for every \(f\in\mathcal H_d(K)\),
\[
\#\left\{P\in(\overline K)^2:
[K(P):K]\leq B\ \text{and}\
\hat h_f(P)<\epsilon h(f)+C
\right\}\leq N.
\]
In particular, the same \(\epsilon,C,N\) must work for every map in the indicated family. The height inequality is strict, whereas the field-degree bound is inclusive. The supplied record labels this assertion open; that label is not an independently verified mathematical assessment.
\subsection{Short English statement}
For fixed \(K,B,d\), the conjecture asks for a uniform bound on the number of algebraic points of field degree at most \(B\) whose height \(\hat h_f\) is less than \(\epsilon h(f)+C\), as the generalized H\'{e}non map \(f\) varies. The constants may depend on \(K,B,d\), but not on \(f\). The map family and height conventions remain subject to source clarification.
\subsection{Sources}
\begin{itemize}
\item Supplied UnsolvedMath local registry record: permanent ID \texttt{801098}, exact source problem code \texttt{catalogue entry}, title \emph{Henon Maps: Open Problems --- Conjecture 50}. The record attributes the problem to Berger and Bedford, with proposed year 2023, and supplies the status \emph{open}.
\item Berger and Bedford, \emph{Henon Maps: A List of Open Problems} (2023), Conjecture 50. Supplied source URL: \url{https://arxiv.org/abs/2312.03907}. The provenance records access on 2026-10-16 and extraction by \texttt{arxiv\_latex}; these are supplied provenance statements, not a claim of access or independent verification here.
\end{itemize}
\end{document}
